\documentclass[a4paper]{article}
\usepackage[pdftex]{graphicx}
\usepackage[utf8]{inputenc}
\usepackage{enumerate}
\usepackage{siunitx}
\sisetup{locale=DE}
\usepackage{href-ul}
\hypersetup{
	colorlinks=true,
	linkcolor=blue,
	urlcolor=blue}
\usepackage{icomma}
\usepackage{amssymb}
\usepackage{geometry}
\geometry{a4paper, top=15mm, left=15mm, right=15mm, bottom=15mm,
	headsep=10mm, footskip=12mm}

\begin{document}
	%Title
	{\bf Impromptu task from physics }\\
	
	\noindent A ball of mass $m=\SI{12}{\milli\gram}$ and charge $Q=\SI{-1,4e-9}{\coulomb}$ is at rest on a horizontal surface and is released at time $t_0=\SI{0}{\second}$ in the middle of the homogeneous field of a plate capacitor. The electric field strength is
	$\SI{36}{\kilo\volt\over \meter}$. Friction phenomena can be neglected.
	
	\begin{enumerate}[1.]
		
		\item Explain why there is no constant charging current when charging the capacitor.
		
		\item Show by general calculation that the speed at which the bead hits the capacitor plate depends on the plate distance $d$:
		$$ v(d)=\SI{2,0}{\sqrt{\meter}\over \second}\cdot\sqrt{d}$$
		
		\item Calculate the voltage $U$ that must be applied to the capacitor so that the bead hits the capacitor plate with the speed $v=\SI{0.56}{\meter\over \second}$.
		
		\item Does this impact speed change if the distance between the plates is increased while there is a connection to the voltage source? Justify your answer.
		
		\item What happens to the bead after it hits the capacitor plate?
		
	\end{enumerate}
	
	\fbox{
		\begin{minipage}{0.5\textwidth}
			For the solution, please \href{https://www.okuyakl.de/phys/p12efeldL093/le093.pdf}{click here} or scan the QR code.\\
			Additional worksheets are available at
			
			\href{http://www.okuyakl.com}{www.okuyakl.com}
		\end{minipage}
		\hfill
		\begin{minipage}{0.4\textwidth}
			\includegraphics[width=1.5 cm]{../../viecher/zwe03}
			\includegraphics[width=3 cm]{qre093}
			\includegraphics[width=2 cm]{../../viecher/afanticon1}
			
	\end{minipage}}
\end{document}